Este informe estudia el comportamiento práctico de tres paradigmas de diseño de
algoritmos sobre un mismo problema: encontrar en una grilla de $N \times M$ el
camino de máximo puntaje desde $(1,1)$ hasta $(N,M)$, con movimientos solo
hacia el este y hacia el sur, sorteando murallas intransitables y minas que
cuestan una vida y restan $P$ puntos. Como el rover dispone de $V$ unidades de
blindaje, el problema combina la búsqueda de un camino con una restricción de
recurso, del mismo modo que la capacidad en el problema de la mochila.

Los paradigmas comparados son la fuerza bruta, implementada como backtracking
recursivo que poda únicamente los nodos de fracaso; la programación dinámica,
que obtiene el óptimo exacto tabulando los $N \cdot M \cdot V$ estados; y una
heurística greedy cuya regla local viene fijada por el enunciado. El objetivo
es contrastar sus curvas empíricas de tiempo y memoria contra las cotas
asintóticas que les corresponden, y cuantificar el costo de la miopía greedy
mediante su brecha respecto al óptimo y su tasa de falla.
